\src{189}
following, for example, from the commutativity of the following diagram of
canonical endomorphisms:
\begin{equation}\tag{7.6}
\resizebox{\textwidth}{!}{$
\left\{
\begin{array}{ccccc}
\operatorname{Ext}_{\Osh_X}^{n-m}(X;\Omega_X^{m},\Omega_X^{n})
&\longleftarrow&
\operatorname{Ext}_{\Osh_X}^{n-m}(X;\Omega_Y^{m},\Omega_X^{n})
&\simeq&
\operatorname{Hom}_{\Osh_X}(X;\Omega_Y^{m},\Omega_Y^{m})\\[3pt]
\Big\downarrow&&&&\Big\downarrow\\[3pt]
\operatorname{Ext}_{\Osh_X}^{n-m}(X;\Osh_X,\Omega_X^{n-m})
&\longleftarrow&
\operatorname{Ext}_{\Osh_X}^{n-m}(X;\Osh_Y,\Omega_X^{n-m})
&\simeq&
\operatorname{Hom}_{\Osh_X}(X;\Omega_X^{m},\Omega_Y^{m})
\end{array}
\right.
$}
\end{equation}

This gives an exact analogue of the formalism of Poincaré duality on compact
oriented varieties. In particular, Theorem 4 determines the cohomology class
associated with the diagonal of $X\times X$. A well-known argument then gives,
for example, a Lefschetz formula.

\result{THEOREM}{5} Let $f$ be an endomorphism of a nonsingular projective
variety $X$ whose fixed points all have multiplicity 1. Then the number of
these fixed points is congruent, modulo the characteristic of $k$, to the
alternating sum of the traces of the endomorphisms of the $H^{i}(X)$ induced by
$f$.

The restriction imposed on $f$ is due to the difficulties mentioned in the
remark in Section 4. Note, however, that the Lefschetz formula remains true if
$f$ is ``homotopic'' to an endomorphism all of whose fixed points have
multiplicity 1.

\fsection{8}{Generalization of the duality theorem}

Let $X$ be a nonsingular algebraic variety such that every coherent algebraic
sheaf $F$ on $X$ is isomorphic to a quotient of a locally free coherent
algebraic sheaf (as is the case if $X$ is locally closed in projective space).
Then every coherent algebraic sheaf $F$ on $X$ admits a finite resolution $L$
by locally free sheaves; for any two such resolutions one can always find a
third and a homomorphism from it to each of the first two compatible with the
augmentations. Likewise, if $L$ is a finite locally free resolution of $F$ and
a homomorphism $F'\longrightarrow F$ is given, there exists a finite locally
free resolution $L'$ of $F'$ and a homomorphism $L'\longrightarrow L$
compatible with $F'\longrightarrow F$, which may be assumed surjective if
$F'\longrightarrow F$ is surjective. Given two integers $r,s\geq0$, this makes
it possible to define two cohomological multifunctors in the arguments
$A_1,\ldots,A_r$; $B_1,\ldots,B_s$,

